\documentclass[../main.tex]{subfiles}
\begin{document}

\subsection{О чём эта глава}%==================================================

До сих пор мы имели дело с движением вдоль прямой и с переходом
между системами отсчёта, движущимися поступательно. Теперь мы
сделаем следующий шаг: рассмотрим движение точки по произвольной
криволинейной траектории.

Примеры таких траекторий встречаются в самых разных задачах.
Камень, брошенный под углом $\alpha$ к горизонту с начальной скоростью
$v_0$, вблизи поверхности Земли летит по параболе:
\[
y = x\tg\alpha - \frac{g\,x^2}{2v_0^2\cos^2\alpha}.
\]

Планеты движутся вокруг Солнца по эллипсам, близким к окружностям:
$\frac{x^2}{a^2} + \frac{y^2}{b^2} = 1$; некоторые кометы пролетают
вблизи Солнца по гиперболическим траекториям: $\frac{x^2}{a^2} - \frac{y^2}{b^2} = 1$.
При сложении двух взаимно перпендикулярных гармонических колебаний
разной частоты траектория точки представляет собой одну из фигур
Лиссажу:
\[
x = A\sin(\omega_1 t + \delta), \qquad y = B\sin\omega_2 t.
\]
Точка на ободе катящегося колеса описывает циклоиду. Если колесо
радиуса $R$ катится без проскальзывания вдоль оси $Ox$, а его центр
движется со скоростью~$v$, то координаты этой точки можно записать
в параметрическом виде:
\[
x = vt - R\sin\frac{vt}{R}, \qquad
y = R\left(1 - \cos\frac{vt}{R}\right).
\]
Если же колесо вращается вокруг неподвижной оси, то та же точка
движется уже по окружности.

Прямолинейное движение можно рассматривать как предельный случай
движения по окружности~--- с бесконечно большим радиусом.
Иными словами, криволинейное движение~--- это более общий класс
движений, чем прямолинейное.

Для любого криволинейного движения существенно, что вектор скорости~$\mbfv$ непрерывно \textit{меняет направление}, а ускорение, вообще
говоря, направлено под углом к скорости. Особенно удобно увидеть это
на примере движения по окружности~--- самого простого криволинейного
движения. Поэтому начнём именно с него.

\subsection{Угловая скорость. Центростремительное ускорение}
%==============================================================================
Сначала введём понятие угловой скорости радиус-вектора в общем случае.
Оно имеет смысл для любого движения точки относительно выбранного
начала координат, не только для движения по окружности.

\begin{definition}
\textit{Угловой скоростью радиус-вектора} называется быстрота изменения
угла поворота $\varphi$ радиус-вектора точки:
\[
\omega = \lim_{\Delta t \to 0}\frac{\Delta\varphi}{\Delta t}\equiv\dv{\varphi}{t} \equiv \dot{\varphi}(t).
\]
\end{definition}

Далее мы опускаем знак предела $\lim$ при работе с малыми величинами и последующим предельным переходом, фактически имея в виду дифференциалы функций. 

Рассмотрим движение точки по окружности радиуса $R$ с центром в начале координат со скоростью $v$. За малое время $\Delta t$ точка проходит дугу длины
$\Delta s \hm = v\,\Delta t$, а её радиус-вектор поворачивается на угол  $\Delta\varphi = \frac{\Delta s}{R} = \frac{v\,\Delta t}{R}$. Разделим это равенство на $\Delta t$, получим:
\[
\omega = \frac{v}{R}.
\]
Отметим, что эта формула справедлива как для равномерного, так и для неравномерного движения по окружности. При этом надо понимать, что центр окружности, по которой движется точка, должен совпадать с началом координат. В противном случае для угловой скорости вращения радиус-вектора получится другое выражение.

Далее в этом разделе для простоты будем рассматривать движение по окружности с постоянной по модулю скоростью.

\begin{definition}
\textit{Периодом} $T$ называется время одного полного оборота.

\textit{Частотой} $\nu$ называется число оборотов в единицу времени.
Для равномерного движения по окружности
\[
T = \frac{2\pi}{\omega} = \frac{2\pi R}{v},
\qquad
\nu = \frac{1}{T} = \frac{\omega}{2\pi}.
\]
\end{definition}

Пусть точка движется по окружности радиусом $R$ с центром в начале координат с постоянной угловой скоростью $\omega$. Будем считать, что в момент $t = 0$ координаты точки равны $(R, 0)$. Тогда в момент времени $t$ её координаты будут равны
\[
x(t) = R\cos(\omega t), \qquad y(t) = R\sin(\omega t).
\]

Найдём теперь проекции скорости по определению. Рассмотрим отношение малого изменения координаты к малому изменению времени. 

Для координаты $x$ имеем
\[
\Delta x = x(t + \Delta t) - x(t)
= R\cos(\omega t + \omega\Delta t) - R\cos(\omega t).
\]
Воспользуемся приближённой формулой $\cos(\varphi + \delta) \approx \cos\varphi - \delta\sin\varphi$, справедливой при малом $\delta$. После преобразований получим для приращения координаты формулу
\[
\Delta x = -R\omega\Delta t \cdot \sin(\omega t).
\]
Аналогично для малого изменения координаты $y$ получаем:
\[
\Delta y = R\omega\Delta t \cdot\cos(\omega t).
\]
Таким образом, компоненты скорости равны:
\[
v_x = -R\omega\sin(\omega t), \qquad v_y =R\omega\cos(\omega t).
\]
Теперь легко проверить, что модуль скорости остаётся постоянным:
\[
v^2 = v_x^2 + v_y^2 = R^2\omega^2.
\]
Кроме того, вычислив скалярное произведение:
\[
\qty(\mbfv\cdot\mbfr) = (-R\omega\sin\omega t)(R\cos\omega t)
+ (R\omega\cos\omega t)(R\sin\omega t) = 0,
\]
убеждаемся в том, что в каждый момент времени вектор скорости~$\mbfv$  \textit{ортогонален} (направлен перпендикулярно) радиус-вектору~$\mbfr$.

При равномерном движении по окружности модуль скорости не меняется;
меняется только её направление. Для того, чтобы найти ускорение, нужно понять, как за малое время меняется сам вектор скорости.

Поступим точно так же: рассмотрим отношение малого изменения проекции скорости к малому изменению времени. Применяя те же приближённые формулы для синуса и косинуса, что и выше, получаем
\[
\Delta v_x = -R\omega^2\Delta t \cdot \cos(\omega t), \qquad
\Delta v_y = -R\omega^2\Delta t \cdot \sin(\omega t),
\]
откуда находим для компонент ускорения:
\[
a_x = -R\omega^2\cos(\omega t), \qquad a_y = -R\omega^2\sin(\omega t).
\]
Сравнив эти выражения с координатами точки, замечаем, что
\[
\mbfa = -\omega^2\mbfr.
\]
Ускорение в каждый момент времени направлено противоположно радиус-вектору~$\mbfr$, то есть к центру окружности; поэтому его называют \textit{центростремительным}. Его модуль постоянен:
\[
a = R\omega^2 = \frac{v^2}{R}.
\]

Покажем и другой вывод этой формулы. 
Рассмотрим  векторы $\mbfv(t)$ и $\mbfv(t + \Delta t)$. Они равны по модулю $v$, но повёрнуты друг относительно друга на угол
$\Delta\varphi$. Если отложить их из одной точки, то концы этих
векторов окажутся на окружности радиуса $v$, а вектор
$\Delta\mbfv = \mbfv(t + \Delta t) - \mbfv(t)$ совпадёт с хордой,
стягивающей дугу, соответствующую углу $\Delta\varphi$. Отметим, что угол между векторами $\Delta\mbfv$ и $\mbfv$ приближённо равен $\frac{\pi}{2}$. С учётом малости угла $\Delta\varphi$ для модуля изменения скорости имеем:
\[
|\Delta\mbfv| = 2v\sin\frac{\Delta\varphi}{2} \approx v \,\Delta\varphi.
\]
Но для движения по окружности $\Delta\varphi = \frac{v\,\Delta t}{R} = \omega\Delta t$, поэтому
\[
a = \frac{|\Delta\mbfv|}{\Delta t}
= \omega v = \frac{v^2}{R} = \omega^2 R.
\]
Поскольку изменение вектора скорости, как уже было сказано, практически ортогонально скорости, вектор ускорения направлен также перпендикулярно скорости, как и следовало ожидать.  Рассмотрим теперь два примера. 

\begin{example}
Материальная точка движется по окружности с постоянной скоростью
$v = 2~\text{м}/\text{с}$. За время $\tau = 2~\text{с}$ вектор скорости повернулся на угол
$\alpha = \frac{\pi}{4}$. Найдите период обращения и центростремительное ускорение точки. При равномерном движении вектор скорости поворачивается с той же угловой
скоростью $\omega$, что и сама точка. Из условия:
$\omega = \frac{\alpha}{\tau} = \frac{\pi/4}{2} = \frac{\pi}{8}~\text{рад}/\text{с}$.

Период обращения:
\[
T = \frac{2\pi}{\omega} = 16~\text{с}.
\]

Центростремительное ускорение:
\[
a = v\omega = 2 \cdot \frac{\pi}{8}
= \frac{\pi}{4}~\text{м}/\text{с}^2 \approx 0{,}785~\text{м}/\text{с}^2.
\]
\end{example}

\begin{example}
Небольшое тело движется по окружности радиуса $R$ с угловой скоростью $\omega$.
Центр окружности совпадает с началом координат; в момент $t = 0$
тело находится на оси $Ox$ с положительной координатой.
Найдите угловую скорость вращения отрезка, соединяющего тело
с точкой $P = (-R,\, 0)$.

Обозначим через $B = (R,\, 0)$ точку окружности, в которой тело
находится в момент $t = 0$. Через время $t$ центральный угол
$\angle AOB$ равен $\omega t$.

Поскольку луч $PB$ направлен вдоль оси $Ox$, угол между отрезком $PA$
и осью $Ox$ равен вписанному углу $\angle APB$. Этот угол опирается на
ту же дугу $AB$, что и центральный угол $\angle AOB$, поэтому
\[
\angle APB = \frac{1}{2}\angle AOB = \frac{\omega t}{2}.
\]
Значит, отрезок $PA$ поворачивается вдвое медленнее, чем радиус-вектор
$OA$. Его угловая скорость равна
\[
\Omega = \frac{\Delta(\omega t/2)}{\Delta t} = \frac{\omega}{2}.
\]
\end{example}

%==============================================================================
\subsection{Нормальное и тангенциальное ускорения}
%==============================================================================

В этом подразделе речь идёт о неравномерном движении по окружности
постоянного радиуса $R$. В этом случае меняются модуль скорости и направление вектора скорости.

Рассмотрим векторы скорости
$\mbfv(t)$ и $\mbfv(t + \Delta t)$. Пусть угол между ними равен~$\Delta\gamma$. Отложим оба вектора из одной точки. Тогда вектор
$\Delta\mbfv = \mbfv(t + \Delta t) - \mbfv(t)$ замыкает этот треугольник скоростей.

Рассмотрим проекции вектора $\mbfv(t + \Delta t)$ на направление 
$\mbfv(t)$ и перпендикулярное направление. Тогда вектор $\Delta\mbfv$
раскладывается на две составляющие: тангенциальную, направленную вдоль~$\mbfv(t)$, и нормальную, перпендикулярную~$\mbfv(t)$. Модули этих составляющих обозначим через $\Delta v_\tau$ и $\Delta v_n$ соответственно. 


Для нормальной составляющей имеем: $\Delta v_n = |\mbfv(t + \Delta t)|\sin\Delta\gamma$. При малом $\Delta t$ угол $\Delta\gamma$ мал, а разность
$|\mbfv(t + \Delta t)| - |\mbfv(t)|$ имеет порядок $v \qty(\Delta \gamma)^{2}$.
Поэтому с точностью до величин второго порядка малости: $\Delta v_n \approx v\,\Delta\gamma$. Разделив на $\Delta t$, получаем
\[
a_n = \frac{\Delta v_n}{\Delta t}
= v\,\frac{\Delta\gamma}{\Delta t}
= v\,\omega_v,
\]
где $\omega_v$ --- угловая скорость вращения вектора скорости. При движении по окружности угол поворота вектора скорости равен углу
поворота радиус-вектора, то есть $\Delta\gamma = \Delta\varphi$, а следовательно $\omega_{v} = \omega$. Таким образом
\[
a_n = v\omega = \frac{v^2}{R}.
\]

Для тангенциальной составляющей: $\Delta v_\tau = |\mbfv(t + \Delta t)|\cos\Delta\gamma - |\mbfv(t)|$. Так как при малом $\Delta\gamma$ выполняется $\cos\Delta\gamma \approx 1$,
то
\[
\Delta v_\tau \approx |\mbfv(t + \Delta t)| - |\mbfv(t)| = \Delta v,
\]
где $\Delta v$ --- изменение модуля скорости. Поэтому
$a_\tau = \frac{\Delta v_\tau}{\Delta t} = \frac{\Delta v}{\Delta t}$.

Итак, тангенциальное ускорение характеризует быстроту изменения модуля скорости, а нормальное --- быстроту изменения направления вектора скорости.

\begin{definition}
При неравномерном движении по окружности постоянного радиуса ускорение
раскладывается на две взаимно перпендикулярные составляющие:
\[
a_n = v\,\omega_v = v\omega = \frac{v^2}{R},
\qquad
a_\tau = \frac{\Delta v}{\Delta t}.
\]
\end{definition}

Рассмотрим связь тангенциального ускорения с быстротой изменения угловой скорости. При движении по окружности постоянного радиуса $\Delta v = R\,\Delta\omega$. Разделив это равенство на $\Delta t$, получаем:
\[
a_\tau = R\,\frac{\Delta\omega}{\Delta t}.
\]
Быстроту изменения угловой скорости называют \textit{угловым ускорением}:
\[
\varepsilon = \frac{\Delta\omega}{\Delta t}.
\]
Тогда
\[
a_\tau = R\varepsilon.
\]
При разгоне $\varepsilon > 0$ и $a_\tau > 0$; при торможении знаки отрицательны.

Если угловое ускорение постоянно, то
формулы для $\omega(t)$ и $\varphi(t)$ совпадают с хорошо знакомыми формулами
равноускоренного движения по прямой, в которых $v$ заменено на $\omega$,
а $a$ --- на $\varepsilon$:
\[
\omega(t) = \omega_0 + \varepsilon t,
\qquad
\varphi(t) = \varphi_0 + \omega_0 t + \frac{\varepsilon t^2}{2}.
\]

Пусть $\theta$ --- острый угол между полным ускорением и тангенциальной
составляющей ускорения. Так как тангенциальная и нормальная составляющие
взаимно перпендикулярны, то
\[
\tg\theta = \frac{a_n}{|a_\tau|}.
\]
В случае разгона, когда $a_\tau > 0$, отсюда получаем
\[
\tg\theta = \frac{a_n}{a_\tau}
= \frac{v^2/R}{R\varepsilon}
= \frac{\omega^2}{\varepsilon}.
\]

\begin{example}
Частица начинает двигаться по окружности радиуса $R$ из состояния
покоя с постоянным тангенциальным ускорением $a_\tau$.
Найдите угол $\theta$ между вектором ускорения и вектором скорости
после первого оборота.

Так как $a_\tau = R\varepsilon$, угловое ускорение $\varepsilon$ постоянно.
Поскольку в начальный момент $\omega_0 = 0$, имеем
\[
\omega = \varepsilon t,
\qquad
\varphi = \frac{\varepsilon t^2}{2}.
\]
После первого оборота $\varphi = 2\pi$, поэтому
\[
2\pi = \frac{\varepsilon t^2}{2}.
\]
Умножая это равенство на $2\varepsilon$, получаем
\[
\omega^2 = \varepsilon^2 t^2 = 4\pi\varepsilon.
\]
Следовательно,
\[
\tg\theta = \frac{\omega^2}{\varepsilon} = 4\pi,
\qquad
\theta = \arctg(4\pi) \approx 85{,}5^\circ.
\]
Так как при разгоне тангенциальное ускорение сонаправлено со скоростью,
этот же угол является углом между полным ускорением и вектором скорости.
\end{example}

Важно, что формулы
\[
a_n = v\,\omega_v,
\qquad
a_\tau = \frac{\Delta v}{\Delta t}
\]
справедливы не только для движения по окружности постоянного радиуса,
но и для любого криволинейного движения. Первая формула связывает
нормальное ускорение с угловой скоростью вращения вектора скорости,
а вторая показывает, что тангенциальное ускорение равно быстроте
изменения модуля скорости.

Особенность движения по окружности постоянного радиуса состоит лишь в том, что здесь
$\omega_v = \frac{v}{R}$, и поэтому нормальное ускорение выражается через
радиус окружности формулой $a_n = \frac{v^2}{R}$. Для произвольной
траектории такая связь требует отдельного геометрического рассмотрения.
Именно к этому мы перейдём в следующем подразделе.

\begin{example}
В момент времени, когда скорость частицы равна $v$, её ускорение равно $a$
по модулю и направлено под углом $\alpha$ к вектору скорости. Какова
угловая скорость $\omega_v$ вращения вектора скорости в этот момент?

Нормальная составляющая ускорения равна
\[
a_n = a\sin\alpha,
\]
так как именно нормальная составляющая перпендикулярна вектору скорости. По формуле
$a_n = v\omega_v$ получаем:
\[
\omega_v = \frac{a_n}{v} = \frac{a\sin\alpha}{v}.
\]
\end{example}

%==============================================================================
\subsection{Радиус кривизны}
%==============================================================================

Малый участок любой криволинейной траектории можно приблизить окружностью.
Однако окружность можно подобрать по-разному, и потому возникает вопрос:
какая из них даёт наилучшее приближение в данной точке?

Введём понятие \textit{соприкасающейся окружности}. Как известно, через любые три точки, не лежащие на одной прямой, проходит единственная окружность. Возьмём на траектории точку $P$ и две близкие к ней точки $Q$ и $S$, расположенные по разные стороны от $P$. Проведём через них окружность. Если теперь неограниченно приближать точки $Q$ и $S$ к точке $P$ вдоль траектории, то эта окружность будет стремиться к некоторой предельной окружности. Такая предельная окружность называется \textit{соприкасающейся}.

Чем замечательна эта окружность? Во-первых, в точке $P$ она касается траектории, то есть имеет с ней общую касательную. Во-вторых, при движении вдоль кривой касательная поворачивается; у соприкасающейся окружности она поворачивается с той же быстротой, что и у самой траектории. Тем самым соприкасающаяся окружность задаёт не только касательную к кривой в точке $P$, но и её локальную кривизну.

\begin{definition}
\textit{Радиусом кривизны} $\rho$ траектории в данной точке называется
радиус соприкасающейся окружности в этой точке.
\end{definition}

Чем сильнее изогнута траектория, тем меньше радиус кривизны. Для прямой
линии $\rho = \infty$.

Теперь свяжем радиус кривизны с нормальным ускорением. Поскольку вблизи точки $P$ соприкасающаяся окружность и траектория имеют
одинаковую быстроту изменения угла поворота касательной, движение по данной
траектории в этой точке можно рассматривать как движение по окружности
радиуса $\rho$. Поэтому для нормального ускорения остаётся справедливой
формула
\[
a_n = v\,\omega_v,
\]
а для соприкасающейся окружности
\[
\omega_v = \frac{v}{\rho}.
\]
Подставляя это выражение в предыдущую формулу, получаем:
\[
a_n = \frac{v^2}{\rho}.
\]
Следовательно,
\[
\rho = \frac{v^2}{a_n}.
\]

Эта формула особенно удобна в задачах: если в некоторый момент известны
скорость точки и нормальное ускорение, то радиус кривизны траектории в
этот момент находится сразу.

Важно понимать, что радиус кривизны --- это геометрическая характеристика траектории. Вообще говоря, радиус кривизны не зависит от того, с какой скоростью точка
движется по траектории. Он определяется формой самой кривой. По одной и той же
траектории точка может двигаться по разным законам, но радиус её кривизны
в данной геометрической точке останется тем же.

Для знакомых с понятием производной приведём формулы, позволяющие находить $\rho$
непосредственно по уравнению кривой. Если траектория задана уравнением
$y = f(x)$, то
\[
\rho = \frac{\bigl(1 + (y')^2\bigr)^{3/2}}{|y''|}.
\]
Если же она задана параметрически: $x = x(t)$, $y = y(t)$, то
\[
\rho = \frac{\bigl(\dot x^2 + \dot y^2\bigr)^{3/2}}{|\dot x\,\ddot y - \dot y\,\ddot x|}.
\]
Здесь параметр $t$ не обязан совпадать со временем.

\begin{example}
Найдём радиус кривизны баллистической траектории в произвольный момент
времени. Пусть тело брошено со скоростью $v_0$ под углом $\alpha$ к
горизонту.

Ускорение тела равно $\mbfg$. Нормальная составляющая этого ускорения равна
\[
a_n = g\cos\theta,
\]
где $\theta$ --- угол между вектором скорости и горизонтом в рассматриваемый
момент. Горизонтальная проекция скорости в процессе движения не меняется, поэтому:
\[
v\cos\theta = v_0\cos\alpha,
\qquad
v = \frac{v_0\cos\alpha}{\cos\theta}.
\]
Отсюда имеем формулу для радиуса кривизны
\[
\rho = \frac{v^2}{a_n}
= \frac{\bigl(v_0\cos\alpha / \cos\theta\bigr)^2}{g\cos\theta}
= \frac{v_0^2\cos^2\alpha}{g\cos^3\theta}.
\]

В верхней точке траектории $\theta = 0$, поэтому
\[
\rho = \frac{v_0^2\cos^2\alpha}{g}.
\]
\end{example}

\begin{example}
Колесо радиуса $R$ катится без проскальзывания по горизонтали со
скоростью центра $v$. На спице, на расстоянии $r < R$ от центра,
закреплена точка. Найдите радиус кривизны её траектории в тот момент,
когда точка находится в верхнем положении.

В этот момент скорость точки складывается из переносной скорости центра
и вращательной скорости относительно центра. Обе они направлены
горизонтально вправо, поэтому
\[
v_{\text{т}} = v + \omega r = v + \frac{v}{R}r = v\,\frac{R+r}{R}.
\]
Ускорение точки в этот момент обусловлено только вращением вокруг центра
колеса, так как центр движется равномерно. Поэтому оно направлено
вертикально вниз и равно
\[
a_n = \omega^2 r = \frac{v^2}{R^2}\,r.
\]
Поскольку скорость точки в этот момент горизонтальна, а ускорение
вертикально, всё ускорение является нормальным. Следовательно,
\[
\rho = \frac{v_{\text{т}}^2}{a_n}
= \frac{v^2 (R+r)^2/R^2}{v^2 r/R^2}
= \frac{(R+r)^2}{r}.
\]
При $r = R$ получаем
\[
\rho = 4R,
\]
что соответствует вершине циклоиды.
\end{example}

%==============================================================================
\subsection{Задачи преследования}
%==============================================================================

В задачах преследования одно тело («преследователь») всегда движется
в направлении другого («цели»). Скорость при этом постоянна по величине,
но непрерывно меняет направление. Ключевым инструментом служит
производная расстояния между преследователем и целью.

Найдём скорость сближения. Пусть $A$~--- преследователь, $B$~--- цель. Обозначим $r = |AB|$.
Скорость изменения расстояния $r$ равна проекции относительной скорости
$\mbfv_B - \mbfv_A$ на прямую $AB$. Поэтому за малое время $\Delta t$
\[
\Delta r = \hat{\mbfe}_{AB}\cdot(\mbfv_B - \mbfv_A)\,\Delta t,
\]
где $\hat{\mbfe}_{AB}$~--- единичный вектор от $A$ к $B$.
Если $\Delta r < 0$, расстояние убывает (тела сближаются).
Раскроем явно: пусть $\beta_A$~--- угол между $\mbfv_A$ и направлением $AB$,
а $\beta_B$~--- угол между $\mbfv_B$ и тем же направлением. Тогда:
\[
\frac{\Delta r}{\Delta t} = v_B\cos\beta_B - v_A\cos\beta_A.
\]

\begin{example}
Четыре черепахи находятся в вершинах квадрата со стороной~$a$. Они
начинают двигаться одновременно, каждая --- со скоростью $v$ к своей
ближайшей соседке по часовой стрелке. Через какое время черепахи
встретятся? Найдите ускорение любой из черепах и радиус кривизны её
траектории в момент начала движения и спустя время $\frac{T}{2}$.

Пусть в некоторый момент времени черепахи находятся в точках
$B_1$, $B_2$, $B_3$ и $B_4$. Благодаря симметрии начального расположения
и одинаковым законам движения, поворот системы на $90^\circ$ вокруг
центра $O$ исходного квадрата циклически переводит текущие положения
черепах друг в друга: $B_1 \to B_2 \to B_3 \to B_4 \to B_1$. Отсюда
следует, что расстояния от черепах до точки $O$ равны
($OB_1 = OB_2 = OB_3 = OB_4$), а угол между направлениями на соседних
черепах составляет $90^\circ$. Таким образом, в любой момент времени
точки $B_1$, $B_2$, $B_3$ и $B_4$ являются вершинами квадрата с
неизменным центром в точке $O$. Из-за монотонного уменьшения стороны
этого квадрата в процессе движения черепахи встретятся в его центре.

Найдём скорость сближения соседних черепах. Пусть $A$ и $B$ --- соседние
черепахи, $r = |AB|$ --- расстояние между ними. По условию черепаха $A$
ползёт прямо к $B$, поэтому $\beta_A = 0$ и $\cos\beta_A = 1$. Черепаха
$B$ ползёт к следующей черепахе под прямым углом к прямой $AB$ (угол
между сторонами квадрата равен $90^\circ$), поэтому $\beta_B = 90^\circ$
и $\cos\beta_B = 0$. По формуле
\[
\frac{\Delta r}{\Delta t} = v_B\cos\beta_B - v_A\cos\beta_A
= v\cdot 0 - v\cdot 1 = -v.
\]
Знак минус означает, что расстояние убывает; скорость сближения равна
$v$ и постоянна. Поэтому время встречи
\[
T = \frac{a}{v}.
\]

Перейдём к ускорению. Скорость каждой черепахи постоянна по модулю,
значит, ускорение чисто нормальное: $w = v\omega$, где $\omega$ ---
угловая скорость вращения вектора скорости. Вектор скорости всегда
направлен на соседнюю черепаху, поэтому $\omega$ совпадает с угловой
скоростью вращения отрезка, соединяющего соседних черепах.

Найдём её. Перейдём в поступательно движущуюся систему отсчёта черепахи
$A$. В этой системе черепаха $B$ движется перпендикулярно прямой $AB$
со скоростью $v$ (поперечная компонента; продольная равна нулю). За
малое время $\Delta t$ отрезок $AB$ поворачивается на угол
\[
\Delta\gamma = \frac{v\,\Delta t}{r},
\]
так что
\[
\omega = \frac{v}{r}.
\]

В начальный момент $r = a$, поэтому
\[
w_0 = v\omega = \frac{v^2}{a},
\qquad
\rho_0 = \frac{v^2}{w_0} = a.
\]

Поскольку скорость сближения постоянна, расстояние между соседними
черепахами убывает равномерно: $r(t) = a - vt$. Через время
$\frac{T}{2} = \frac{a}{2v}$ имеем $r = \frac{a}{2}$, поэтому
\[
w_1 = \frac{v^2}{\frac{a}{2}} = \frac{2v^2}{a},
\qquad
\rho_1 = \frac{v^2}{w_1} = \frac{a}{2}.
\]
\end{example}
%
%\figtikz{images/kin/cur_turtles.tikz}{Три черепахи в вершинах правильного треугольника: каждая бежит к соседке. Скорость сближения соседних черепах равна $\tfrac{3}{2}v$.}{fig:cur_turtles}

\begin{example}
На берегу находятся пункты $A$ и $B$, удалённые на расстояние $L$ друг от
друга. Контрабандисты отправляются из пункта $A$ в открытое море
перпендикулярно берегу с постоянной скоростью $v$. Офицер береговой охраны,
обнаружив нарушителей в момент отчаливания их судна, тотчас пускается
в погоню на катере из точки $B$. Катер всё время держит курс на
нарушителей и, двигаясь с постоянной скоростью $u$, настигает судно на
расстоянии $L$ от берега. Во сколько раз $k = \frac{u}{v}$ скорость катера
превосходит скорость судна?

Пусть $r$ --- расстояние между катером и судном, $h$ --- разность их
расстояний до берега, $\varphi$ --- угол между вертикалью и прямой,
соединяющей катер с судном. За малое время $\Delta t$ судно удаляется
от берега на $v\,\Delta t$, а катер проходит $u\,\Delta t$ вдоль этой
прямой.

Спроектируем скорости на прямую, соединяющую катер с судном. Изменение
расстояния между ними равно
\[
\Delta r = v\cos\varphi\,\Delta t - u\,\Delta t.
\]
Спроектировав скорости на вертикаль, найдём изменение величины $h$:
\[
\Delta h = v\,\Delta t - u\cos\varphi\,\Delta t.
\]

Умножим первое равенство на $u$, второе --- на $v$ и сложим. Слагаемые,
содержащие $\cos\varphi$, сокращаются, и остаётся
\[
u\,\Delta r + v\,\Delta h = (v^2 - u^2)\,\Delta t.
\]
Поскольку $u$ и $v$ постоянны, левая часть есть изменение величины
$ur + vh$ за время $\Delta t$.

Просуммируем это равенство по всем малым промежуткам от начала погони
до поимки. Слева получится полное изменение величины $ur + vh$, а справа
--- $(v^2 - u^2)T$, где $T$ --- время погони:
\[
\bigl(ur + vh\bigr)\Big|_{\text{нач}}^{\text{кон}} = (v^2 - u^2)T.
\]
В начале погони катер находится в $B$, а судно --- в $A$, поэтому
$r = L$, $h = 0$. В момент поимки $r = 0$ и $h = 0$. Следовательно,
\[
(0 + 0) - (uL + 0) = (v^2 - u^2)T,
\]
то есть
\[
uL = (u^2 - v^2)T.
\]

Судно удалилось от берега на $L$ со скоростью $v$, поэтому $T = \frac{L}{v}$.
Подставляя $u = kv$, получаем
\[
kvL = (k^2v^2 - v^2)\,\frac{L}{v} = vL(k^2 - 1).
\]
Сокращая на $vL$, приходим к уравнению
\[
k = k^2 - 1
\quad\Rightarrow\quad
k^2 - k - 1 = 0.
\]
С учётом $k > 1$ выбираем положительный корень:
\[
k = \frac{1 + \sqrt{5}}{2} \approx 1{,}618.
\]
Это число --- золотое сечение.
\end{example}

\end{document}